% GENERATED FILE. Edit canonical lessons, then run npm run generate:books.
% canonical-source-hash: fnv1a64:721b0ae916104be3
% canonical-lessons: PT-C221-genro, PT-C221-nora, PT-C221-noivo, PT-C221-noiva, PT-C221-viuvo-a-viuva

\chapter{Son-in-law, Daughter-in-law, Fiancé, Fiancée, Widower}
\label{ch:pt-son-in-law-daughter-in-law-fiance-fiancee-widower}
\hlchaptermodality{221}

\begin{chapteropening}
\textbf{By the end of this chapter:} \emph{I can say a son-in-law, a daughter-in-law, a fiancé, a fiancée and a widower.}

Complete the last lesson of chapter 221: I can say a son-in-law, a daughter-in-law, a fiancé, a fiancée and a widower.
\end{chapteropening}

\section[o genro]{o genro — a son-in-law}

\label{lesson:PT-C221-genro}



\begin{quote}
Before the new one: say the Portuguese for a pharmacist, then the Portuguese for a butcher.
\end{quote}

\subsection*{You'll want to know: o genro}
\textbf{o genro} — \textquotedblleft{}a son-in-law\textquotedblright{}.

\begin{grammarlens}[title={What you've built}]
One more word for this chapter.
\end{grammarlens}

\subsection*{Your turn}
\begin{itemize}
\raggedright
  \item \emph{Say it:} \emph{o genro}
  \item \emph{Say it:} \emph{o genro}, once more
  \item \emph{Say it:} say \emph{o genro}
  \item \emph{Recall:} say the Portuguese for a pharmacist, then the Portuguese for a butcher, then say \emph{o genro} again
\end{itemize}

\begin{tcolorbox}[breakable,colback=teal!4,colframe=teal!35!black,title={Before you move on}]
What does o genro mean? (\textquotedblleft{}A son-in-law\textquotedblright{}.) Say it once more.
\end{tcolorbox}
\section[a nora]{a nora — a daughter-in-law}

\label{lesson:PT-C221-nora}



\begin{quote}
Before the new one: say the Portuguese for a butcher, then the Portuguese for a son-in-law.
\end{quote}

\subsection*{You'll want to know: a nora}
\textbf{a nora} — \textquotedblleft{}a daughter-in-law\textquotedblright{}.

\begin{grammarlens}[title={What you've built}]
One more word for this chapter.
\end{grammarlens}

\subsection*{Your turn}
\begin{itemize}
\raggedright
  \item \emph{Say it:} \emph{a nora}
  \item \emph{Say it:} \emph{a nora}, once more
  \item \emph{Say it:} say \emph{a nora}
  \item \emph{Recall:} say the Portuguese for a butcher, then the Portuguese for a son-in-law, then say \emph{a nora} again
\end{itemize}

\begin{tcolorbox}[breakable,colback=teal!4,colframe=teal!35!black,title={Before you move on}]
What does a nora mean? (\textquotedblleft{}A daughter-in-law\textquotedblright{}.) Say it once more.
\end{tcolorbox}
\section[o noivo]{o noivo — a fiancé, a bridegroom}

\label{lesson:PT-C221-noivo}



\begin{quote}
Before the new one: say the Portuguese for a son-in-law, then the Portuguese for a daughter-in-law.
\end{quote}

\subsection*{You'll want to know: o noivo}
\textbf{o noivo} — \textquotedblleft{}a fiancé, a bridegroom\textquotedblright{}.

\begin{grammarlens}[title={What you've built}]
One more word for this chapter.
\end{grammarlens}

\subsection*{Your turn}
\begin{itemize}
\raggedright
  \item \emph{Say it:} \emph{o noivo}
  \item \emph{Say it:} \emph{o noivo}, once more
  \item \emph{Say it:} say \emph{o noivo}
  \item \emph{Recall:} say the Portuguese for a son-in-law, then the Portuguese for a daughter-in-law, then say \emph{o noivo} again
\end{itemize}

\begin{tcolorbox}[breakable,colback=teal!4,colframe=teal!35!black,title={Before you move on}]
What does o noivo mean? (\textquotedblleft{}A fiancé, a bridegroom\textquotedblright{}.) Say it once more.
\end{tcolorbox}
\section[a noiva]{a noiva — a fiancée, a bride}

\label{lesson:PT-C221-noiva}



\begin{quote}
Before the new one: say the Portuguese for a daughter-in-law, then the Portuguese for a fiancé.
\end{quote}

\subsection*{You'll want to know: a noiva}
\textbf{a noiva} — \textquotedblleft{}a fiancée, a bride\textquotedblright{}.

\begin{grammarlens}[title={What you've built}]
One more word for this chapter.
\end{grammarlens}

\subsection*{Your turn}
\begin{itemize}
\raggedright
  \item \emph{Say it:} \emph{a noiva}
  \item \emph{Say it:} \emph{a noiva}, once more
  \item \emph{Say it:} say \emph{a noiva}
  \item \emph{Recall:} say the Portuguese for a daughter-in-law, then the Portuguese for a fiancé, then say \emph{a noiva} again
\end{itemize}

\begin{tcolorbox}[breakable,colback=teal!4,colframe=teal!35!black,title={Before you move on}]
What does a noiva mean? (\textquotedblleft{}A fiancée, a bride\textquotedblright{}.) Say it once more.
\end{tcolorbox}
\section[o viúvo / a viúva]{o viúvo / a viúva — a widower, a widow}

\label{lesson:PT-C221-viuvo-a-viuva}



\begin{quote}
Before the new one: say the Portuguese for a fiancé, then the Portuguese for a fiancée.
\end{quote}

\subsection*{You'll want to know: o viúvo / a viúva}
\textbf{o viúvo / a viúva} — \textquotedblleft{}a widower, a widow\textquotedblright{}.

\begin{grammarlens}[title={What you've built}]
One more word for this chapter.
\end{grammarlens}

\subsection*{Your turn}
\begin{itemize}
\raggedright
  \item \emph{Say it:} \emph{o viúvo / a viúva}
  \item \emph{Say it:} \emph{o viúvo / a viúva}, once more
  \item \emph{Say it:} say \emph{o viúvo / a viúva}
  \item \emph{Recall:} say the Portuguese for a fiancé, then the Portuguese for a fiancée, then say \emph{o viúvo / a viúva} again
\end{itemize}

\begin{tcolorbox}[breakable,colback=teal!4,colframe=teal!35!black,title={Before you move on}]
What does o viúvo / a viúva mean? (\textquotedblleft{}A widower, a widow\textquotedblright{}.) Say it once more.
\end{tcolorbox}
